\chapter{公共场所卫生}
\setcounter{chapter}{10}
\chapterintro{
了解公共场所的概念、分类及卫生学特点；掌握公共场所的基本卫生要求；了解公共场所的卫生管理与监督内容。
}

\section{公共场所概述}

\begin{defbox}
\textbf{公共场所（Public Place）}：根据公众生活活动和社会活动的需要，人工建成的具有多种服务功能的封闭式和开放式的公共建筑设施。
\end{defbox}

\subsection{卫生学特点}

\begin{enumerate}[leftmargin=*]
    \item \imp{人群密集，流动性大}——易造成疾病传播
    \item \imp{设备及物品易被污染}——公共用品反复使用
    \item \imp{涉及面广}——各类人群均可能接触
    \item \imp{从业人员素质参差不齐}——管理水平差异大
\end{enumerate}

\subsection{公共场所的分类}

\begin{multicols}{2}
\begin{enumerate}[leftmargin=*]
    \item 住宿与交际场所（宾馆等）
    \item 洗浴与美容场所
    \item 文化娱乐场所
    \item 体育与游乐场所
    \item 文化交流场所（图书馆等）
    \item 购物场所（商场等）
    \item 就诊与交通场所
\end{enumerate}
\end{multicols}

\section{公共场所的卫生要求与监督}

\textbf{基本卫生要求：}空气质量达标、微小气候适宜、采光照明良好、噪声符合标准、公用物品清洁、卫生设施齐全。

\textbf{经营企业自身卫生管理：}配备卫生管理人员和制度；从业人员卫生培训；健康检查；对顾客卫生宣教。

\textbf{卫生机构管理：}从业人员培训体检；发放卫生许可证；健康教育；经常性卫生监督。

\begin{defbox}
\textbf{经常性卫生监督}：对公共场所日常卫生状况进行的定期巡回监督检查，确保持续符合卫生要求。
\end{defbox}

\section{复习题}

\begin{enumerate}[leftmargin=*, itemsep=8pt]
    \item \textbf{【名词解释】}公共场所；经常性卫生监督
    \item \textbf{【简答题】}简述公共场所的卫生学特点。
    \item \textbf{【简答题】}简述公共场所经营企业自身卫生管理的主要内容。
\end{enumerate}

\subsection*{参考答案要点}

\begin{enumerate}[leftmargin=*, itemsep=5pt]
    \item \textbf{公共场所}：根据公众生活和社会活动需要人工建成的具有多种服务功能的公共建筑设施。\textbf{经常性卫生监督}：对公共场所日常卫生状况的定期巡回监督检查。

    \item 四大特点：人群密集流动性大、设备物品易被污染、涉及面广、从业人员素质参差不齐。

    \item 配备卫生管理人员和制度；组织从业人员学习卫生知识技能；组织健康检查；开展对顾客卫生宣教。
\end{enumerate}

\newpage

% =====================================================================
%  CHAPTER 10: 环境质量评价
% =====================================================================
\newpage